\section{Introduction}
\label{sec:introduction}

Atmospheric retrieval is the inverse problem of inferring atmospheric and
planetary properties from a measured spectrum. A forward model maps a parameter
vector $\boldsymbol{\theta}$ to a spectrum $\boldsymbol{m}(\boldsymbol{\theta})$;
an inference method then evaluates which regions of parameter space are
consistent with the observation. The method has developed from terrestrial and
Solar-system remote sensing, including optimal estimation and the NEMESIS
framework \citep{Rodgers2000,Irwin2008}, into a central tool for interpreting
transiting exoplanets \citep[e.g.][]{MadhusudhanSeager2009,Line2013}.

Modern retrieval packages make different trade-offs between physical breadth,
speed, configurability, and software transparency. TauREx provides a modular
Python retrieval architecture and a plugin system \citep{AlRefaie2021};
petitRADTRANS combines line-by-line and correlated-$k$ radiative transfer with
retrieval tooling \citep{Molliere2019,Molliere2020}; POSEIDON emphasizes
multidimensional atmospheres and broad observing geometries
\citep{MacDonald2017,MacDonald2024}; and PICASO couples reflected-light,
thermal-emission, transit, climate, and scattering calculations
\citep{Batalha2019,Batalha2022}. CHIMERA, Brewster, ARCiS, PLATON, and
NEMESIS/NemesisPy span complementary use cases from hot-Jupiter transmission
and emission to directly imaged planets and brown dwarfs
\citep{Line2013,Burningham2017,Min2020,Zhang2019,Irwin2008}.

No single framework is uniformly optimal. Broad packages can make many science
cases accessible, but tightly coupled data loading, physics, and inference can
make controlled replacement or validation of one component difficult.
Research-oriented codes may be fast and scientifically productive while
depending on project-specific scripts, implicit conventions, or external data
layouts. Conversely, strict abstraction can obscure the numerical path or
introduce complexity before the underlying physics is validated. These are
engineering trade-offs rather than a ranking of scientific quality.

JWST has made these trade-offs especially consequential. Its precision and
broad wavelength coverage expose molecular features, cloud effects, and
longitudinal temperature structure simultaneously. For example, the
$2$--$12\,\mu$m emission spectrum of WASP-69b is inconsistent with a simple
homogeneous clear atmosphere and favours models containing aerosols and a
dayside temperature gradient \citep{Schlawin2024}. WASP-80b provides a cooler
comparison target: NIRCam transmission and emission spectra yielded strong
methane evidence \citep{Bell2023}, and its later panchromatic emission spectrum
constrained metallicity and C/O over $2.4$--$12\,\mu$m
\citep{WASP80Panchromatic2025}.

Here we introduce \ROBERT{} (\roberttodo{expand the acronym, or state explicitly
that ROBERT is a name rather than an acronym}), a Python package built around
explicit scientific data models and independently testable numerical
components. Its initial science focus is JWST thermal emission, while its
interfaces avoid encoding an emission-only assumption. The distribution is
named \code{robert-exoplanets} and is currently pre-release software.

This paper has two aims. First, we describe the architecture, physics, inference
interfaces, and validation strategy of \ROBERT{}. Secondly, we use a reanalysis
of WASP-80b to test whether a reproducible end-to-end implementation recovers
the published atmospheric constraints and whether additional temperature,
cloud, chemistry, or multidataset flexibility is supported by the data. Section
\ref{sec:methods} describes the model and benchmarks, Section
\ref{sec:results} defines and will report the WASP-80b analysis, and Section
\ref{sec:conclusions} summarizes the present scope and future development.

